\chapter{Sequential Linear Programming -- Implementation Details}
\label{sec:SLP_appendix}

This appendix expands the Sequential Linear Programming (SLP)
description given in \autoref{sec:SLP}. The purpose is to
document the numerical procedure used to optimize the LT1292,
including the linearized subproblem, slack relaxation, merit
function, filter acceptance, sufficient-decrease test, adaptive
move-limit trust region and stopping criteria.

The implemented optimizer is a Python port of the merit-function
variant of the MATLAB framework \texttt{fminslp} described by
\cite{sorensen2026fminslp}. The port reproduces every algorithmic
ingredient of the reference: a slack-relaxed LP subproblem, a
linear--quadratic merit penalty with a fixed slack weight, a
Fletcher--Leyffer--Toint global convergence filter, a
sufficient-decrease test on the merit function and adaptive
component-wise move limits driven by the iterate history. The
solver is coupled to the MAPDL analysis pipeline described in the
code-implementation chapter. At each outer iteration, the current
structural design is evaluated in MAPDL, finite-difference
sensitivities are computed, and a linear programming (LP)
subproblem is solved inside a move-limit trust region. The
resulting trial design is then re-evaluated in MAPDL before being
accepted or rejected.

The appendix is structured as follows. \autoref{sec:slpA:problem}
defines the optimization problem and the sign convention used by the
implementation. \autoref{sec:slpA:linearise} describes the local
linearization and the finite-difference gradient calculation.
\autoref{sec:slpA:lp} presents the slack-relaxed LP subproblem.
\autoref{sec:slpA:merit} introduces the merit function used both
in the LP objective and in the nonlinear acceptance test.
\autoref{sec:slpA:accept} explains the filter, sufficient-decrease
test and trust-region backtracking. \autoref{sec:slpA:movelim}
describes the adaptive move limits, which act as a component-wise
trust region. \autoref{sec:slpA:termination} describes the
termination criteria. The numerical settings used in the present
study are listed in \autoref{sec:slpA:settings}.

% =====================================================================
\section{Problem statement and sign convention}
\label{sec:slpA:problem}

The structural optimization problem is formulated as a constrained
mass-minimization problem. The design vector is denoted
$x \in \mathbb{R}^n$, where each component represents a physical
design variable in the MAPDL beam model, such as a tube radius,
diameter or wall thickness. The problem is written as
\begin{equation}
\label{eq:appA_problem}
\begin{aligned}
    \min_{x \in \mathbb{R}^n} \quad & f(x), \\
    \text{s.t.} \quad
    & c_i(x) \geq 0,
    \qquad i = 1,\dots,m, \\
    & x^L \leq x \leq x^U .
\end{aligned}
\end{equation}
Here, $f(x)$ is the structural mass returned by the MAPDL analysis.
The vector $c(x)$ collects the design constraints, including
utilization constraints, cross-section class constraints,
eigenvalue constraints and segment-mass constraints. The lower and
upper bounds $x^L$ and $x^U$ define the physically admissible design
space.

In the thesis formulation, the constraints are written in the
non-negative convention $c_i(x) \geq 0$. In this convention, a
positive value means that the constraint is satisfied, while a
negative value means that the constraint is violated. The SLP
implementation uses the opposite convention,
\begin{equation}
\label{eq:appA_g_convention}
    g_i(x) \leq 0,
    \qquad i = 1,\dots,m .
\end{equation}
The two formulations are connected by a sign change in the constraint
callback,
\begin{equation}
\label{eq:appA_sign_flip}
    g(x) = -c(x).
\end{equation}
All algorithmic quantities in the following sections are therefore
defined using the $g(x) \leq 0$ convention. Thus, a constraint is
satisfied when $g_i(x) \leq 0$ and violated when $g_i(x)>0$.

% =====================================================================
\section{Local linearization and finite-difference gradients}
\label{sec:slpA:linearise}

Sequential Linear Programming solves a nonlinear constrained problem
by repeatedly replacing it with a local linear approximation. At
iteration $k$, the current design is $x_k$, and the trial design is
written as
\begin{equation}
\label{eq:appA_trial_design}
    x_{\text{trial}} = x_k + \Delta x ,
\end{equation}
where $\Delta x$ is the design change proposed by the LP subproblem.

The objective function is approximated by the first-order Taylor
expansion
\begin{equation}
\label{eq:appA_lin_f}
    f(x_k + \Delta x)
    \approx
    f(x_k) + \nabla f(x_k)^{\!\top}\Delta x .
\end{equation}
Similarly, the constraint vector is approximated as
\begin{equation}
\label{eq:appA_lin_g}
    g(x_k + \Delta x)
    \approx
    g(x_k) + J(x_k)\Delta x .
\end{equation}
The Jacobian $J(x_k) \in \mathbb{R}^{m \times n}$ contains the local
constraint sensitivities,
\begin{equation}
\label{eq:appA_jacobian}
    J_{ij}(x_k)
    =
    \frac{\partial g_i(x_k)}{\partial x_j},
    \qquad
    i = 1,\dots,m,\quad j = 1,\dots,n .
\end{equation}

The MAPDL model is treated as a black-box analysis model. Analytical
gradients are therefore not available. Instead, the objective
gradient and the constraint Jacobian are estimated by forward finite
differences. For design variable $i$, the perturbation size is chosen
as
\begin{equation}
\label{eq:appA_fd_step}
    h_i =
    h_{\mathrm{rel}}
    \max\left(1, |x_{k,i}|\right),
\end{equation}
where $h_{\mathrm{rel}}$ is a fixed relative finite-difference step.
Using the unit vector $e_i$, the objective sensitivity is estimated
as
\begin{equation}
\label{eq:appA_fd_f}
    \frac{\partial f(x_k)}{\partial x_i}
    \approx
    \frac{
        f(x_k + h_i e_i) - f(x_k)
    }{h_i},
\end{equation}
and the corresponding constraint sensitivities are estimated as
\begin{equation}
\label{eq:appA_fd_g}
    \frac{\partial g_j(x_k)}{\partial x_i}
    \approx
    \frac{
        g_j(x_k + h_i e_i) - g_j(x_k)
    }{h_i},
    \qquad j = 1,\dots,m .
\end{equation}

For $n$ design variables, one gradient evaluation requires one
baseline MAPDL evaluation and $n$ perturbed evaluations. The finite
difference calculation is therefore the dominant computational cost
of each outer SLP iteration.

% =====================================================================
\section{Slack-relaxed linear subproblem}
\label{sec:slpA:lp}

If the local linear models in \autoref{eq:appA_lin_f} and
\autoref{eq:appA_lin_g} are used directly, the SLP subproblem becomes
\begin{equation}
\label{eq:appA_direct_lp}
\begin{aligned}
    \min_{\Delta x}\quad
        & \nabla f(x_k)^{\!\top}\Delta x ,
    \\
    \text{s.t.}\quad
        & g(x_k) + J(x_k)\Delta x \leq 0 ,
    \\
        & x^{L,\mathrm{cur}}
          \leq
          x_k + \Delta x
          \leq
          x^{U,\mathrm{cur}} .
\end{aligned}
\end{equation}
The constant term $f(x_k)$ is omitted from the objective because it
does not affect the minimizer of the LP.

The direct LP in \autoref{eq:appA_direct_lp} may become infeasible if
the current design violates one or more constraints and the
move-limit box is too small to recover feasibility in one step. To
make the LP solvable from any iterate, each constraint is relaxed by
a non-negative slack variable $y_i$,
\begin{equation}
\label{eq:appA_slack}
    g_i(x) - y_i \leq 0,
    \qquad
    y_i \geq 0,
    \qquad
    i = 1,\dots,m .
\end{equation}
The slack variable represents the amount of artificial constraint
relaxation. If $y_i = 0$, the original constraint is enforced. If
$y_i > 0$, the corresponding constraint is temporarily relaxed.
Because permanent slack usage would correspond to an infeasible
engineering design, the slacks are penalized through the merit
function introduced in \autoref{sec:slpA:merit}.

The initial slack vector is chosen as
\begin{equation}
\label{eq:appA_slack_init}
    y_i^0 = \max\left(0, g_i(x_0)\right),
    \qquad i = 1,\dots,m ,
\end{equation}
so that the initial slack values cover the initial constraint
violations. To keep the LP bounded, the slacks are also assigned a
large finite upper bound,
\begin{equation}
\label{eq:appA_slack_upper}
    y_i^U =
    10^6 \max\left(1,\max_j g_j(x_0)\right),
    \qquad i = 1,\dots,m .
\end{equation}
This upper bound is chosen large enough that it does not normally
restrict the solution, but it prevents the LP from becoming
unbounded in pathological cases.

The augmented variable and its increment are defined as
\begin{equation}
\label{eq:appA_augmented_variable}
    z =
    \begin{bmatrix}
        x \\ y
    \end{bmatrix},
    \qquad
    \Delta z =
    \begin{bmatrix}
        \Delta x \\ \Delta y
    \end{bmatrix}.
\end{equation}
The slack-relaxed LP solved at each outer iteration is then
\begin{equation}
\label{eq:appA_slack_relaxed_lp}
\begin{aligned}
    \min_{\Delta z}\quad
        & \nabla_z \Phi(x_k,y_k)^{\!\top}\Delta z ,
    \\
    \text{s.t.}\quad
        & g(x_k) + J(x_k)\Delta x - (y_k + \Delta y) \leq 0 ,
    \\
        & x^{L,\mathrm{cur}}
          \leq
          x_k + \Delta x
          \leq
          x^{U,\mathrm{cur}},
    \\
        & 0
          \leq
          y_k + \Delta y
          \leq
          y^U .
\end{aligned}
\end{equation}
The current bounds $x^{L,\mathrm{cur}}$ and $x^{U,\mathrm{cur}}$ are
the move-limit bounds described in \autoref{sec:slpA:movelim}, and
the merit-gradient cost vector $\nabla_z\Phi$ is defined in
\autoref{sec:slpA:merit}.

The linearized relaxed constraint in \autoref{eq:appA_slack_relaxed_lp}
can be written in the standard LP form
\begin{equation}
\label{eq:appA_lp_standard_form}
    A_{\mathrm{ub}}\Delta z \leq b_{\mathrm{ub}},
\end{equation}
where
\begin{equation}
\label{eq:appA_lp_standard_matrices}
    A_{\mathrm{ub}}
    =
    \begin{bmatrix}
        J(x_k) & -I_m
    \end{bmatrix},
    \qquad
    b_{\mathrm{ub}}
    =
    -\left(g(x_k)-y_k\right).
\end{equation}
Because each inequality has its own non-negative slack variable with
a large finite upper bound, the LP remains practically feasible for
the constraint violations encountered in the present study. The LP is
solved by the HiGHS dual-simplex routine.

% =====================================================================
\section{Merit function with linear-quadratic slack penalty}
\label{sec:slpA:merit}

The slack variables make the LP robust against infeasible iterates,
but they must not remain active in the final design. The algorithm
therefore uses a scalar merit function that combines the objective
value and the slack violation into one quantity. The merit function
is defined as
\begin{equation}
\label{eq:appA_merit}
    \Phi(x,y)
    =
    f(x)
    +
    a_F
    \sum_{i=1}^{m}
    \left(
        R y_i
        +
        \frac{1}{2}y_i^2
    \right),
\end{equation}
where $R > 0$ is the slack penalty and
\begin{equation}
\label{eq:appA_af}
    a_F = \max\left(|f(x_0)|,1\right)
\end{equation}
is an objective scaling factor fixed at the first MAPDL evaluation.
In contrast to many SLP variants, the slack penalty $R$ is held
\emph{constant} throughout the run: it is the user-supplied
infeasibility penalization parameter and never grows or shrinks. The
slack term in \autoref{eq:appA_merit} therefore always combines a
linear part $R y_i$ that makes even small slack values expensive with
a quadratic part $\tfrac{1}{2}y_i^2$ that increasingly penalizes
larger violations. The scaling factor $a_F$ makes the slack penalty
less sensitive to the units and magnitude of the mass objective.

The merit function in \autoref{eq:appA_merit} has two roles. First,
it discourages the optimizer from using slacks permanently. Second,
it provides a scalar quantity that can be used to compare the
predicted LP improvement with the actual nonlinear improvement after
a trial design has been evaluated in MAPDL.

The gradient of the merit function with respect to the augmented
variable $z=[x,y]$ is
\begin{equation}
\label{eq:appA_merit_grad}
    \nabla_z \Phi(x,y)
    =
    \begin{bmatrix}
        \nabla f(x)
        \\
        a_F(R\mathbf{1}+y)
    \end{bmatrix},
\end{equation}
where $\mathbf{1}$ is a vector of ones of length $m$. This gradient is
used as the cost vector in the LP objective in
\autoref{eq:appA_slack_relaxed_lp}. The LP therefore seeks a step that
reduces the linearized merit function, not only the structural mass.
This couples the LP subproblem to the nonlinear acceptance test
described in \autoref{sec:slpA:accept}.

The implementation also exposes an augmented-Lagrangian variant in
which the slack penalty in \autoref{eq:appA_merit} is replaced by
$a_F\sum_i\bigl(y_i\lambda_i+\tfrac{1}{2}Ry_i^2\bigr)$ and the
multiplier vector $\lambda$ is updated after every accepted step by
$\lambda \leftarrow \max(\lambda + Ry, 1)$. The merit-function
variant in \autoref{eq:appA_merit} is used in the present study; the
LP subproblem, the filter and the trust-region machinery are
identical for both variants.

% =====================================================================
\section{Step acceptance}
\label{sec:slpA:accept}

The LP solution provides a candidate step
$\Delta z=[\Delta x,\Delta y]$, but the LP itself does not guarantee
that the step is good for the original nonlinear MAPDL model. This is
because the LP is based only on local linear approximations. The
trial design must therefore be evaluated in MAPDL and accepted only
if it gives sufficient progress.

The implementation combines two acceptance mechanisms: a global
convergence filter following \cite{sorensen2015thicknessfilters} and a
sufficient-decrease test on the merit function. A trial point is
accepted only if it passes both. Trial points that fail either test
are rejected, the move-limit trust region is contracted, and the LP
is re-solved.

% ---------------------------------------------------------------------
\paragraph{Filter acceptance.}

The filter compares trial points using two quantities: a measure of
relaxed constraint violation and a scaled merit value. The relaxed
maximum constraint violation is defined as
\begin{equation}
\label{eq:appA_filter_h}
    h(x,y)
    =
    \max_i
    \max\left(0, g_i(x)-y_i\right),
\end{equation}
and the scaled merit value is
\begin{equation}
\label{eq:appA_filter_f}
    f_{\mathrm{filter}}(x,y)
    =
    \frac{\Phi(x,y)}{\Phi_0},
    \qquad
    \Phi_0=\max\!\left(|\Phi(x_0,y_0)|,10^{-30}\right) .
\end{equation}
The small lower bound on $\Phi_0$ is a numerical safeguard against
division by zero. A trial point is represented by the pair
\begin{equation}
\label{eq:appA_filter_pair}
    \left(h^*, f^*\right)
    =
    \left(
        h(x_{\mathrm{trial}},y_{\mathrm{trial}}),
        f_{\mathrm{filter}}(x_{\mathrm{trial}},y_{\mathrm{trial}})
    \right).
\end{equation}

The filter $\mathcal{F}_k$ accepts the trial point if it is not
dominated by any stored filter entry. The implemented test is
\begin{equation}
\label{eq:appA_filter_accept}
    h^* \leq \beta h_j
    \quad \text{or} \quad
    f^* + \gamma h^* \leq f_j,
    \qquad
    \forall\,(h_j,f_j)\in\mathcal{F}_k ,
\end{equation}
with filter slackness parameters $\beta = 1-10^{-6}$ and
$\gamma = 10^{-6}$. In words, the trial point must give either
sufficiently lower relaxed violation or sufficiently lower scaled
merit value compared with every stored filter entry. The filter is
initialised with the single permissive entry $(\infty,\infty)$, so
the first trial point is always filter-acceptable regardless of its
feasibility. After an accepted infeasible step ($h^*>0$), the new
filter pair is inserted into the filter and entries dominated by the
new pair are removed. Accepted points with $h^*=0$ are not stored, so
the filter never grows when the optimization remains feasible with
respect to the relaxed inequalities.

% ---------------------------------------------------------------------
\paragraph{Sufficient decrease in the merit function.}

Filter acceptance alone is not sufficient, because the filter does
not directly compare the actual nonlinear reduction with the
reduction predicted by the LP model. The implementation therefore
also checks a merit-function decrease condition.

The actual merit decrease at the trial point is
\begin{equation}
\label{eq:appA_merit_actual}
    \Delta\Phi_{\mathrm{act}}
    =
    \Phi(x_k,y_k)
    -
    \Phi(x_k+\Delta x, y_k+\Delta y),
\end{equation}
and the predicted merit decrease from the LP model is
\begin{equation}
\label{eq:appA_merit_pred}
    \Delta\Phi_{\mathrm{pred}}
    =
    -\nabla_z\Phi(x_k,y_k)^{\!\top}\Delta z .
\end{equation}
A filter-acceptable trial point is then accepted by the
sufficient-decrease condition
\begin{equation}
\label{eq:appA_armijo}
    \Delta\Phi_{\mathrm{act}} \geq \sigma\Delta\Phi_{\mathrm{pred}},
    \qquad
    \sigma = 2\cdot 10^{-6} ,
\end{equation}
whenever the LP predicts a strict merit decrease,
$\Delta\Phi_{\mathrm{pred}}>0$. If $\Delta\Phi_{\mathrm{pred}}\leq 0$,
the LP step is not a predicted descent step for the merit function;
the sufficient-decrease test is then not informative, and the trial
point is accepted on the basis of the filter alone.

% ---------------------------------------------------------------------
\paragraph{Trust-region backtracking.}

If the filter rejects the trial point, or if the filter accepts but
the sufficient-decrease condition in \autoref{eq:appA_armijo} fails
(with $\Delta\Phi_{\mathrm{pred}}>0$), the trial step is rejected.
The move-limit box is then contracted using the reduction rule
described in \autoref{sec:slpA:movelim}, and the LP is re-solved with
the smaller bounds. This is not a fractional line search along the
previous step direction, but a re-solution of the LP inside a smaller
component-wise trust region. Each inner attempt counts towards the
function-evaluation budget of \autoref{sec:slpA:termination}, and the
inner loop continues until either a trial point is accepted or one of
the convergence indicators in \autoref{sec:slpA:termination} fires.

% =====================================================================
\section{Adaptive move limits as a trust region}
\label{sec:slpA:movelim}

The local linear approximation is only reliable near the current
design. The SLP method therefore restricts each LP step to a
component-wise move-limit box,
\begin{equation}
\label{eq:appA_move_box}
    x_i^{L,\mathrm{cur}}
    \leq
    x_{k,i} + \Delta x_i
    \leq
    x_i^{U,\mathrm{cur}}, \qquad i = 1,\dots,n .
\end{equation}
This box acts as a component-wise trust region. Large steps are
avoided because they would make the linearized objective and
constraints unreliable.

The move limits are updated separately for each design variable. This
is important for the LT1292 because the design variables have
different physical scales. For example, chord diameters may be of
order $60~\mathrm{mm}$, while brace wall thicknesses may be of order
$2~\mathrm{mm}$. A single scalar move limit in physical units would
therefore be inappropriate.

The trust region for variable $i$ is represented by a non-negative
half-width $\delta_{i,k}$. At iteration $k$, the inherited
half-width is read directly from the previous outer LP bounds,
\begin{equation}
\label{eq:appA_prev_halfwidth}
    \delta_{i,k}^{\mathrm{prev}}
    =
    \tfrac{1}{2}
    \left(
        x_i^{U,\mathrm{cur}} - x_i^{L,\mathrm{cur}}
    \right) .
\end{equation}
The update of $\delta_{i,k}$ is driven by the ratio between the
two most recent accepted design changes,
\begin{equation}
\label{eq:appA_sign_ratio}
    s_{i,k}
    =
    \frac{x_{i,k-1}-x_{i,k-2}}{x_{i,k}-x_{i,k-1}} .
\end{equation}
If $|x_{i,k}-x_{i,k-1}|>\epsilon_{\mathrm{hist}}$ with
$\epsilon_{\mathrm{hist}}=10^{-10}$, the half-width is updated by
\begin{equation}
\label{eq:appA_halfwidth_update}
    \delta_{i,k}
    =
    \begin{cases}
        a_{\mathrm{shr}}\,\delta_{i,k}^{\mathrm{prev}},
            & s_{i,k} < 0 \quad \text{(oscillatory),}
        \\[2pt]
        a_{\mathrm{exp}}\,\delta_{i,k}^{\mathrm{prev}},
            & s_{i,k} \geq 0 \quad \text{(coherent),}
    \end{cases}
\end{equation}
with shrink factor $0<a_{\mathrm{shr}}<1$ and expansion factor
$a_{\mathrm{exp}}>1$. Note that $\operatorname{sign}(s_{i,k})$ equals
the sign of the product $(x_{i,k-1}-x_{i,k-2})(x_{i,k}-x_{i,k-1})$, so
a negative ratio means that variable $i$ has changed direction
between iterations and a non-negative ratio means that it has moved
in the same direction.

If the previous accepted change is numerically negligible,
$|x_{i,k}-x_{i,k-1}|\leq\epsilon_{\mathrm{hist}}$, the
iterate history is not considered reliable and the half-width is
reset using the base move-limit fraction $\delta_0$,
\begin{equation}
\label{eq:appA_halfwidth_reset}
    \delta_{i,k} = \delta_0\,\delta_{i,k}^{\mathrm{prev}} .
\end{equation}
At the very first outer iteration, the history is initialised with
$x_{k-1}=x_{k-2}=x_0$, so \autoref{eq:appA_halfwidth_reset} fires
unconditionally. With the initial bounds set to
$(x_i^{L,\mathrm{cur}},x_i^{U,\mathrm{cur}})=(x_i^L,x_i^U)$, the first
half-width is therefore $\tfrac{1}{2}\delta_0\,(x_i^U-x_i^L)$.

After the update, $\delta_{i,k}$ is capped at a fraction of the
global design range,
\begin{equation}
\label{eq:appA_halfwidth_cap}
    \delta_{i,k}
    \leq
    \delta_0\,\left(x_i^U-x_i^L\right) ,
\end{equation}
so the half-width can never grow beyond a prescribed fraction of the
physically admissible interval. The current LP bounds are then formed
as
\begin{equation}
\label{eq:appA_current_bounds}
\begin{aligned}
    x_i^{L,\mathrm{cur}}
    &=
    \max\left(x_i^L, x_{k,i}-\delta_{i,k}\right),
    \\
    x_i^{U,\mathrm{cur}}
    &=
    \min\left(x_i^U, x_{k,i}+\delta_{i,k}\right) .
\end{aligned}
\end{equation}
A symmetric minimum-width guard ensures
\begin{equation}
\label{eq:appA_min_box}
    x_i^{U,\mathrm{cur}}-x_i^{L,\mathrm{cur}} \geq \delta_{\min},
    \qquad
    \delta_{\min} = 2\,\varepsilon_x ,
\end{equation}
where $\varepsilon_x$ is the step-length tolerance of
\autoref{sec:slpA:termination}. This guard prevents the box from
collapsing to zero width and triggers the step-length stopping
indicator before any numerical degeneracy can occur.

The same update is reused for the trust-region backtracking described
in \autoref{sec:slpA:accept}. When a trial step is rejected, both
branches of \autoref{eq:appA_halfwidth_update} are forced to use the
shrink factor $a_{\mathrm{shr}}$, so the half-width contracts
regardless of the sign of $s_{i,k}$. The cap in
\autoref{eq:appA_halfwidth_cap} and the floor in
\autoref{eq:appA_min_box} remain in force. The LP is then re-solved
inside the smaller box. Variables that oscillate are therefore
treated more conservatively, while variables that move consistently
are allowed to take larger steps, and any rejection contracts the
trust region without changing the current iterate.

% =====================================================================
\section{Termination criteria}
\label{sec:slpA:termination}

After each LP solution and trial-point evaluation, four convergence
indicators are evaluated regardless of whether the trial point was
filter-accepted. The first is the merit first-order indicator,
restricted to the design variables,
\begin{equation}
\label{eq:appA_stop_grad}
    \left\|\nabla_x\Phi(x_k,y_k)\right\|_2
    \leq
    \varepsilon_g .
\end{equation}
At a feasible iterate with $y_k=0$, the merit gradient reduces to the
objective gradient and \autoref{eq:appA_stop_grad} becomes the
classical first-order stationarity condition
$\|\nabla f(x_k)\|_2\leq\varepsilon_g$.

The second indicator is the merit-change tolerance, evaluated between
the current iterate and the trial point,
\begin{equation}
\label{eq:appA_stop_merit}
    \left|
        \Phi(x_k,y_k) - \Phi(x_k+\Delta x,y_k+\Delta y)
    \right|
    \leq
    \varepsilon_f .
\end{equation}
The third indicator is the trial step length,
\begin{equation}
\label{eq:appA_stop_step}
    \left\|\Delta x\right\|_2
    \leq
    \varepsilon_x ,
\end{equation}
and the fourth is the objective-limit safeguard,
\begin{equation}
\label{eq:appA_stop_obj}
    \Phi(x_k+\Delta x,y_k+\Delta y) \leq \Phi_{\min} ,
\end{equation}
where $\Phi_{\min}$ is a user-supplied lower bound on the merit
value. The four indicators are tested in the order in which they
appear above, and the first satisfied indicator terminates the
optimization successfully. This precedence follows the reference
MATLAB implementation: the indicators are evaluated independently of
feasibility, so the solver may terminate even if a small residual
slack persists at the final iterate. Engineering feasibility of the
returned design is therefore verified \emph{post-hoc} by inspecting
$\max_i\max(0,g_i(x_k))$ and $\max_i y_{i,k}$ in the optimization
log; a positive slack value at termination is treated as a signal
that the local linear model has become unreliable rather than as a
successful design.

The optimization is additionally bounded by two budget guards: the
outer-iteration counter is limited by $N_{\max}$, and the cumulative
number of LP trial-point evaluations is limited by
$N_{\mathrm{fev}}^{\max}$. Either limit terminates the run as a
non-success.

% =====================================================================
\section{Numerical settings}
\label{sec:slpA:settings}

\autoref{tab:slp_settings} lists the numerical values used in the
present study. The worst-case number of MAPDL evaluations per outer
iteration is
\begin{equation}
\label{eq:appA_fea_cost}
    1+n+N_{\mathrm{bt}},
\end{equation}
where the first evaluation is the baseline analysis at $x_k$, the
next $n$ evaluations are the forward finite-difference perturbations
in \autoref{eq:appA_fd_f}--\autoref{eq:appA_fd_g}, and
$N_{\mathrm{bt}}$ is the number of inner trial-point evaluations
performed during trust-region backtracking. The implementation does
not impose an explicit per-iteration cap on $N_{\mathrm{bt}}$; the
global limit $N_{\mathrm{fev}}^{\max}$ in
\autoref{sec:slpA:termination} acts as the budget guard.

\begin{table}[h]
\centering
\small
\begin{tabular}{lll}
\toprule
Parameter & Value & Description \\
\midrule
$\delta_0$                         & $0.10$           & Initial / maximum move-limit fraction \\
$a_{\mathrm{exp}}$                 & $1.1$            & Move-limit expansion factor \\
$a_{\mathrm{shr}}$                 & $0.5$            & Move-limit shrink factor \\
$\epsilon_{\mathrm{hist}}$         & $10^{-10}$       & Iterate-history activity threshold \\
$R$                                & $10^{3}$         & Fixed slack penalty \\
$\beta$                            & $1-10^{-6}$      & Filter violation parameter \\
$\gamma$                           & $10^{-6}$        & Filter merit parameter \\
$\sigma$                           & $2\cdot 10^{-6}$ & Sufficient-decrease parameter \\
$\varepsilon_g$                    & $10^{-6}$        & Optimality tolerance \\
$\varepsilon_f$                    & $10^{-6}$        & Merit-change tolerance \\
$\varepsilon_x$                    & $10^{-8}$        & Step-length tolerance \\
$\Phi_{\min}$                      & $-10^{20}$       & Objective-limit safeguard \\
$N_{\max}$                         & $1000$           & Maximum outer iterations \\
$N_{\mathrm{fev}}^{\max}$          & $1000$           & Maximum LP trial-point evaluations \\
$h_{\mathrm{rel}}$                 & $10^{-3}$        & Relative finite-difference step \\
\bottomrule
\end{tabular}
\caption{Numerical parameters used for the LT1292 optimization.}
\label{tab:slp_settings}
\end{table}
